\section{Artificial Neural Networks}
\label{sec:background_ann}

\subsection{What is an Artificial Neural Network?}

\glspl{ann} are a family of computational models loosely inspired by the structure of biological neural tissue \cite{lecun2015deep}.
An \gls{ann} is composed of layers of interconnected nodes, or neurons, each of which computes a weighted sum of its inputs and passes the result through a non-linear activation function.
By stacking multiple of these layers in sequence, the network learns a hierarchical composition of transformations that progressively maps raw input data to a task-specific output representation.
The dominant activation function in modern deep learning is the \gls{relu}, defined as $f(x) = \max(0, x)$, which is computationally inexpensive and empirically effective at enabling deep networks to learn complex, non-linear functions \cite{lecun2015deep}.

\subsection{Training an Artificial Neural Network}

Training an \gls{ann} is performed by minimising a differentiable loss function over the training set via gradient descent.
Gradients are propagated from the output layer back through the network using the chain rule of calculus, using a procedure known as backpropagation \cite{lecun2015deep}.
A fundamental requirement of this process is that every function in the computational graph must be differentiable: if a non-differentiable operation is encountered, the chain rule breaks down and no gradient signal can propagate through it.
This constraint is unproblematic for standard \glspl{ann} where activations such as \gls{relu} are differentiable almost everywhere, but becomes a significant challenge for the spiking paradigm, which is non-differentiable, as discussed in Section~\ref{sec:background_snn}.
